\documentclass[10pt]{article}
\usepackage[letterpaper,margin=0.72in,columnsep=0.28in]{geometry}
\usepackage{graphicx}
\graphicspath{{results/plots/}}
\usepackage{booktabs}
\usepackage{amsmath}
\usepackage{hyperref}
\usepackage{caption}
\usepackage[T1]{fontenc}
\usepackage{xcolor}
\hypersetup{colorlinks=true, linkcolor=blue!50!black, urlcolor=blue!50!black, citecolor=blue!50!black}
\captionsetup{font=small, labelfont=bf}
\setlength{\parindent}{0pt}
\setlength{\parskip}{0.45em}
\renewcommand{\arraystretch}{1.05}

\newcommand{\code}[1]{\texttt{\small #1}}

\title{\vspace{-1.5em}}
\date{}

\begin{document}

\twocolumn[
  \begin{center}
    {\LARGE \bfseries A From-Scratch Derivatives Pricing \& Risk Engine}\par\vspace{5pt}
    {\large Validating Three Pricing Methods, an Arbitrage-Free SVI Surface, and Three Greek Estimators Against a Real SPY Book}\par\vspace{8pt}
    \textit{Code: \href{https://github.com/yassineerraji/Derivatives-Pricing-Risk-Engine}{GitHub Repository}}\par\vspace{8pt}
    Yassine Erraji\par\vspace{2pt}
    September 2026
  \end{center}
  \vspace{0.9em}

  {\bfseries \noindent Abstract\hspace{0.6em}}
  This report documents an options pricing and risk-management engine built entirely from first principles (NumPy/SciPy only, no pricing library) and validated at every stage against a closed-form baseline and live SPY market data. Three independent pricing methods --- closed-form Black-Scholes, risk-neutral Monte Carlo, and Crank-Nicolson finite differences --- agree with each other to within a few cents, and a control-variate-plus-antithetic Monte Carlo estimator cuts simulation variance $5.7\times$ (189.1 to 33.3) without introducing bias, matching the standard optimal-coefficient result $\mathrm{Var}(Y_c)=\mathrm{Var}(Y)(1-\rho^2)$. Fitting an SVI volatility surface independently per maturity to a live SPY options chain initially violates static no-arbitrage conditions on 5 of 8 maturity slices; adding two nonlinear constraints to the same weighted-least-squares fit (Durrleman's butterfly condition and a calendar-monotonicity condition against the previous slice) recovers a surface passing both checks on all 8 slices, at a small, documented cost in fit quality. Three Greek estimators --- analytical, bump-and-reprice finite differences, and Monte Carlo pathwise/likelihood-ratio --- agree with the analytical values within their respective standard errors, with the likelihood-ratio estimator carrying markedly higher variance than pathwise at the same path count. Finally, simulating a daily-rebalanced delta hedge for an 8-leg SPY book across 200 independent price paths isolates transaction cost from smile-mismatch and discretization effects: frictional hedging P\&L trails frictionless replication by a mean \$1{,}035, a gap that matches simulated transaction cost almost exactly once cash compounding is accounted for.
  \vspace{1.3em}
]

\section{Introduction}
Derivatives desks price and risk-manage options with a mix of closed-form models, Monte Carlo simulation, PDE solvers, and calibrated volatility surfaces, each validated against the others before being trusted with real risk. This project builds that stack from scratch: pricing, calibration, Greeks, and risk models are implemented directly in NumPy/SciPy, with QuantLib used only inside \code{tests/} to cross-validate the implementation, never as a component of it. Black-Scholes closed-form pricing is the validation baseline throughout, not a component to be skipped past --- every other method is judged both by how well and how fast it reproduces it, and by what happens once it is asked to do something Black-Scholes structurally cannot: price off a smile, estimate a Greek by simulation, or hedge under transaction costs. Every table, plot, and calibrated surface below is produced by a single script under \code{scripts/} against a fixed random seed, so the results are exactly reproducible from a clean checkout given network access to Yahoo Finance.

\section{Methodology \& Conventions}
All pricing and Greeks code assumes a complete, frictionless market and prices under the risk-neutral measure $Q$, where every discounted asset is a martingale and the underlying drifts at $r-q$, never the real-world drift. Black-Scholes-Merton further assumes $S_T$ is log-normal with constant $\sigma$ under $Q$ --- an assumption Section~4 shows the market violates directly. Day counts are ACT/365 with continuously-compounded rates and annualized volatility throughout, stated once here rather than re-derived per module. Market data is a single live SPY options chain pulled via \code{yfinance} and cached locally; sparse or missing strikes are dropped, never fabricated. VaR/ES in Section~6 is the one deliberate exception to $Q$: it replays real historical returns under the physical measure $P$, since ``how much could I lose'' and ``what is the fair hedging price'' are different questions that call for different measures.

\section{Pricing Methods}
Table~1 compares closed-form Black-Scholes against risk-neutral Monte Carlo (200{,}000 paths, antithetic + control variate) and Crank-Nicolson finite differences for European calls ($S_0{=}100$, $T{=}1$y, $r{=}3\%$, $q{=}1\%$, $\sigma{=}20\%$). All three methods agree to within a few cents at every strike; finite differences is consistently the more accurate of the two numerical methods at this grid resolution, while Monte Carlo carries a wider but correctly-sized confidence interval.

\begin{table}[h]
\centering \small
\begin{tabular}{@{}rrrr@{}}
\toprule
$K$ & BS price & MC price (err.) & FD price (err.) \\
\midrule
80  & 22.319 & 22.333 (0.0146) & 22.318 (0.0009) \\
100 & 8.827  & 8.842 (0.0143)  & 8.825 (0.0024) \\
120 & 2.522  & 2.548 (0.0262)  & 2.522 (0.0009) \\
\bottomrule
\end{tabular}
\caption{Call price and absolute error vs. Black-Scholes.}
\end{table}

\textbf{Variance reduction.} The Monte Carlo estimator targets $\theta=E[Y]$ using $X{=}S_T$ as a control variate, whose mean $E[S_T]=S_0e^{(r-q)T}$ is known in closed form. For any constant $b$, $Y_c(b)=Y-b(X-E[X])$ is unbiased, and its variance $\mathrm{Var}(Y)-2b\,\mathrm{Cov}(X,Y)+b^2\mathrm{Var}(X)$ is minimized at $b^{*}=\mathrm{Cov}(X,Y)/\mathrm{Var}(X)$, giving $\mathrm{Var}(Y_c(b^{*}))=\mathrm{Var}(Y)(1-\rho^2)\le\mathrm{Var}(Y)$ always. Table~2 confirms this on a live 200k-path ATM SPY call run: the control variate alone cuts variance $5.5\times$ (189.1 to 34.2); combined with antithetic variates it drops further to 33.3, a smaller, largely independent reduction on top. Figure~1 shows both raw and variance-reduced MC error converging at the theoretical $O(1/\sqrt{n})$ rate, with the reduced estimator consistently below the raw one at every path count.

\begin{table}[h]
\centering \small
\begin{tabular}{@{}llrr@{}}
\toprule
Antithetic & Control var. & Variance & Std. error \\
\midrule
No  & No  & 189.14 & 0.0308 \\
No  & Yes & 34.16  & 0.0131 \\
Yes & No  & 55.25  & 0.0235 \\
Yes & Yes & 33.26  & 0.0182 \\
\bottomrule
\end{tabular}
\caption{MC variance, ATM SPY call, 200k paths.}
\end{table}

\begin{figure}[h]
\centering
\includegraphics[width=\linewidth]{mc_convergence.png}
\caption{MC convergence: mean absolute error vs.\ path count, raw vs.\ antithetic+control-variate, against the $O(1/\sqrt{n})$ reference.}
\end{figure}

\section{Volatility Surface Calibration}
Implied vols are recovered per (strike, maturity) via Brent's method, then fit to Gatheral's SVI parameterization per maturity slice by vega-weighted least squares, so illiquid wing quotes count less than liquid near-the-money ones. Figure~2 fits a live 8-maturity SPY chain ($\sim$4 days to $\sim$10 months): implied vol varies substantially by both strike and maturity, a direct falsification of Black-Scholes' constant-$\sigma$ assumption. The smile carries two shapes Black-Scholes cannot produce: a negative skew (OTM puts trade at higher implied vol than OTM calls of equal moneyness, consistent with the leverage effect and priced crash risk) and a term structure incompatible with any single constant $\sigma$.

An unconstrained per-slice fit, however, is not automatically usable for risk. Run independently, 5 of 8 maturity slices failed Durrleman's butterfly no-arbitrage condition ($g(k)\ge0$ for all log-moneyness $k$) and 4 of 7 consecutive maturity pairs failed the calendar condition (total variance non-decreasing in $T$) --- an artifact of fitting noisy quotes slice-by-slice with no penalty for a non-convex or maturity-decreasing curve, not evidence the market itself is arbitrageable. Adding both conditions as nonlinear SLSQP constraints, fit in increasing-$T$ order against each slice's already-fitted predecessor, recovers a surface passing both checks on all 8 slices (Table~3) at a small, documented cost in fit quality --- the surface Section~6's risk book is actually priced off.

\begin{table}[h]
\centering \small
\begin{tabular}{@{}rr@{}}
\toprule
Maturity (yrs) & Butterfly slack $\min_k g(k)$ \\
\midrule
0.011 & 0.0009 \\
0.019 & 0.0584 \\
0.071 & 0.0424 \\
0.115 & 0.0596 \\
0.211 & 0.1125 \\
0.323 & 0.1474 \\
0.537 & 0.1375 \\
0.819 & 0.0340 \\
\bottomrule
\end{tabular}
\caption{All 8 slices pass Durrleman's butterfly condition ($g\ge0$); all 7 consecutive maturity pairs also pass the calendar condition.}
\end{table}

\begin{figure}[h]
\centering
\includegraphics[width=\linewidth]{vol_surface.png}
\caption{SPY implied volatility: market quotes vs.\ the constrained arbitrage-free SVI fit, per maturity slice.}
\end{figure}

\section{Greeks}
Table~4 compares delta and vega for an ATM SPY call across three estimator families: closed-form analytical, bump-and-reprice finite differences, and two Monte Carlo estimators built on the same GBM path generator used in Section~3 --- pathwise and likelihood-ratio. Finite differences reproduces the analytical value to under $10^{-7}$ error at essentially no extra runtime cost over the closed form. Both Monte Carlo estimators are unbiased and land within roughly one to two standard errors of the analytical value at 200k paths, but the likelihood-ratio estimator's standard error is $2$--$4\times$ larger than pathwise's for the same path count (vega: 0.167 vs.\ 0.626) --- the familiar cost of differentiating the density rather than the payoff, paid here as variance rather than bias.

\begin{table}[h]
\centering \small
\begin{tabular}{@{}llrr@{}}
\toprule
Greek & Estimator & Value & Std. error \\
\midrule
Delta & Analytical & 0.5735 & --- \\
Delta & Finite diff. & 0.5735 & --- \\
Delta & MC pathwise & 0.5724 & 0.0013 \\
Delta & MC likelihood-ratio & 0.5785 & 0.0031 \\
Vega  & Analytical & 38.715 & --- \\
Vega  & Finite diff. & 38.715 & --- \\
Vega  & MC pathwise & 38.871 & 0.167 \\
Vega  & MC likelihood-ratio & 40.304 & 0.626 \\
\bottomrule
\end{tabular}
\caption{ATM SPY call Greeks, 200k MC paths.}
\end{table}

\section{Risk Management}
The book is 8 positions on SPY: mixed calls/puts across three strikes and three maturities (3m/6m/9m), long and short, priced off Section~4's calibrated surface rather than one flat $\sigma$ --- each position looks up its own SVI slice (sticky-strike).

\textbf{VaR/ES.} Table~5 reports 95\%/99\% VaR and ES with confidence intervals from two independent sources: historical simulation on 250 days of real SPY returns, and Monte Carlo simulation at the book's vega-weighted implied vol (20{,}000 scenarios). The two broadly agree at 99\% (VaR within 0.5\%) but diverge more at 95\% (\$701 historical vs.\ \$657 MC), consistent with historical returns carrying fatter realized tails than a single-vol lognormal approximation captures.

\begin{table}[h]
\centering \small
\begin{tabular}{@{}llrr@{}}
\toprule
Method & Conf. & VaR (95\% CI) & ES (95\% CI) \\
\midrule
Historical & 95\% & 701 (676, 712) & 711 (702, 715) \\
Historical & 99\% & 716 (712, 717) & 717 (714, 718) \\
Monte Carlo & 95\% & 657 (652, 661) & 692 (689, 694) \\
Monte Carlo & 99\% & 712 (711, 713) & 716 (715, 716) \\
\bottomrule
\end{tabular}
\caption{Book VaR/ES, \$, with confidence intervals.}
\end{table}

\textbf{Hedging P\&L.} Figure~3 simulates a daily-rebalanced delta hedge across 200 independent SPY price paths, comparing frictionless Black-Scholes/SVI replication P\&L against frictional P\&L that additionally pays half the bid-ask spread plus a market-impact cost on every rebalance. The frictionless mean, $+\$5{,}200$ (std \$4{,}094), is \emph{not} itself ``the cost of hedging'': it is dominated by smile-mismatch P\&L (the book's eight legs span 14.1--21.2\% implied vol, but any one simulated path realizes exactly one vol) and, to a lesser extent, discretization noise from rebalancing daily rather than continuously. Both runs trade identical share quantities along identical paths, so transaction cost is the only thing that differs between them: frictional P\&L trails frictionless by a mean \$1{,}035 (final means \$5{,}200 vs.\ \$4{,}165), matching mean simulated transaction cost of \$1{,}053 to within compounding. That isolated gap, not the raw frictionless mean, is this project's answer to ``what does hedging really cost.''

\begin{figure}[h]
\centering
\includegraphics[width=\linewidth]{hedging_pnl.png}
\caption{Cumulative hedging P\&L, frictionless vs.\ frictional, 200 SPY paths (shaded: 10th--90th percentile).}
\end{figure}

\section{Limitations \& Future Work}
Four scope limits are worth naming rather than leaving implicit. The SVI surface is fit per-slice-sequentially, not as one joint parameterization (e.g.\ SSVI); each slice is constrained only against its immediate predecessor, sufficient for the pairwise checks run here but not a by-construction guarantee at the surface level. The hedging simulation is delta-only: it does not vega- or vanna-hedge, so the smile-mismatch P\&L in Section~6 is a real, unmanaged risk in this book, not a modeling gap. Market impact is a standard linear-cost form, not fit to real order-book data. And every result above is for one underlying (SPY) from one live chain pull; the pipeline runs identically on any liquid ticker (exposed directly in the accompanying Streamlit app), but only SPY is reported here.

\section{Conclusion}
Across four independently testable modules, this engine's from-scratch implementations reproduce their closed-form and cross-validated references to the tolerances the build plan set out: sub-cent pricing agreement, a proven and measured variance reduction, an SVI surface made genuinely arbitrage-free rather than merely fit, and Greek estimators consistent with their analytical values. The risk module's contribution is the one number none of the others can produce on their own: a real book's hedging cost, isolated to \$1{,}035, cleanly separated from the smile-mismatch and discretization effects it would otherwise be entangled with.

{\small
\begin{thebibliography}{9}
\bibitem{bs} F. Black and M. Scholes, ``The Pricing of Options and Corporate Liabilities,'' \textit{Journal of Political Economy}, 81(3), 1973.
\bibitem{merton} R. C. Merton, ``Theory of Rational Option Pricing,'' \textit{Bell Journal of Economics and Management Science}, 4(1), 1973.
\bibitem{gatheral} J. Gatheral, ``A Parsimonious Arbitrage-Free Implied Volatility Parameterization with Application to the Valuation of Volatility Derivatives,'' \textit{Global Derivatives \& Risk Management}, 2004.
\bibitem{durrleman} V. Durrleman, ``From Implied to Spot Volatilities,'' PhD thesis, Princeton University, 2004.
\bibitem{glasserman} P. Glasserman, \textit{Monte Carlo Methods in Financial Engineering}. Springer, 2003.
\bibitem{cranknicolson} J. Crank and P. Nicolson, ``A Practical Method for Numerical Evaluation of Solutions of Partial Differential Equations of the Heat-Conduction Type,'' \textit{Mathematical Proceedings of the Cambridge Philosophical Society}, 43(1), 1947.
\end{thebibliography}
}

\end{document}
